\chapter{Civil Society, Citizenship, Ideology, Pressure Groups \& Social Stratification}

\section{Postmodern \& State-Centred Theories of Power}

\subsection{Baudrillard: Simulacra \& Hyperreality}

\begin{KeyConcept}
\textbf{The conflict between real and hyperreal:}
\begin{itemize}
\item \textbf{Jean Baudrillard} (\textit{End of Politics}): what we see through the media is \textbf{fake news}; the modern form of conflict is \textbf{between the real and the unreal/hyperreal} --- not class conflict (hence also a critique of Marx). \textbf{Hyperreal = simulacra}: a copy of something that does not exist in reality, but assumes a reality of its own.
\item Examples: a fake Nike shoe from Palika Bazaar is a \textit{copy of something that exists} (original Nike); but \textbf{Disneyland, fake news, Fanta} (a replica of an ``orange'' that isn't real) are copies of things that do not exist.
\end{itemize}
\end{KeyConcept}

\begin{ExampleBox}
\textbf{The media-saturated, post-truth world:}
\begin{itemize}
\item The \textbf{Presidential Debate} / TV debates between BJP and Congress spokespersons are media \textit{shows} that present a conflict where none really exists.
\item \textbf{Pokemon Go} (Pikachu on screen but not on the road), \textbf{virtual/artificial reality}, \textbf{Instagram filters/TikTok reels}, \textbf{yellow journalism}, the \textbf{post-truth world}; scientists' \textbf{``Oumuamua''}-style interstellar objects described like sci-fi (aliens). Communal conflicts today are often sparked by viral news that was never real.
\item Tomorrow's world will be \textbf{media-saturated, hyper-real politics}.
\end{itemize}
\end{ExampleBox}

\begin{QuickRevision}
\begin{itemize}
\item Baudrillard: real vs hyperreal = the modern conflict; simulacra = copies of things that don't exist; media $\rightarrow$ fake news $\rightarrow$ post-truth.
\item Examples: Fanta, fake Nike, Disneyland, Pokemon Go, TV debates, yellow journalism.
\end{itemize}
\end{QuickRevision}

\sectiondivider

\subsection{Foucault: Disciplinary Power \& the Panopticon}

\begin{KeyConcept}
\textbf{The Panopticon:}
\begin{itemize}
\item \textbf{Panopticon} (pan = all; optic = visibility) = a prison design with a central watch-tower (often empty) that makes prisoners \textit{feel} watched, so they discipline themselves. Foucault: \textbf{every modern institution --- school, hospital, barracks, factory --- is modelled on the prison}; together they form a \textbf{disciplinary society}.
\item A school is a prison: an authoritative structure, a dress code, emphasis on silence and order, negative reinforcement, walking lines, \textbf{loss of individual autonomy} (autonomy is seen as creating disorder), no input in decision-making, enforced/routinised timetables.
\item Foucault gives no normative judgment --- he only shows \textbf{how power operates} in modern society.
\end{itemize}
\end{KeyConcept}

\begin{ExampleBox}
\textbf{Disciplinary power at scale:}
\begin{itemize}
\item \textbf{USA as the watch-tower}: the \textbf{NPT (Nuclear Proliferation Treaty)}, 1974, divided the world into nuclear-weapon states (the P5) and non-nuclear-weapon states; India is not a signatory (Nehru gave the UNSC seat to China), so India's nuclear programme (Pokhran, Operation Smiling Buddha) had to escape American satellite surveillance --- \textbf{disciplinary punishment}, not physical.
\item \textbf{The census enumerator} collects household data (how many kids, when married), which lets the state regulate you (population control).
\item \textbf{Sweden's ``fat tax''} (data on who eats junk food); \textbf{China's Social Credit System} (2020, universalised; CCTV everywhere; you lose points for VPN/Facebook/smoking and are deprived of services) --- disciplinary power in full play.
\end{itemize}
\end{ExampleBox}

\begin{QuickRevision}
\begin{itemize}
\item Panopticon $\rightarrow$ every modern institution (school, hospital, barracks, factory) modelled on prison $\rightarrow$ disciplinary society.
\item Examples: NPT/Pokhran (USA surveillance), census/fat-tax/social-credit (state regulates through data).
\end{itemize}
\end{QuickRevision}

\sectiondivider

\subsection{Foucault: Power/Knowledge}

\begin{KeyConcept}
\textbf{Power is not possessed:}
\begin{itemize}
\item Unlike other theories (Marxist, pluralist, elite, power-elite, feminist --- all \textit{locate} power in a group), Foucault says \textbf{power cannot be located}. \textbf{Power and knowledge are intertwined}: power is a function of \textbf{knowledge}; knowledge is produced by \textbf{discourse}; discourse is produced by \textbf{specialists} (doctors, professors, administrators).
\end{itemize}
\end{KeyConcept}

\begin{ExampleBox}
\textbf{Madness and sexuality:}
\begin{itemize}
\item The insane were once \textbf{worshipped} (magicians/prophets); homosexuality and transgenders were \textbf{normal} in ancient India (Ardhanarishvara, the Vedas, Khajuraho). With modernity came \textbf{hospitals}, and new disciplines (\textbf{psychology, sexology}) that \textit{classified} some people as abnormal --- filtering out the ``unproductive'' for industrial capitalism. Our knowledge of normal/abnormal is \textbf{manufactured} by specialists and normalised through discourse.
\item We consume this knowledge (via psychology, sexology, bureaucracy) and start treating people as normal/abnormal, legal/illegal, moral/immoral --- everyday practising social exclusion and inequality. This is where power lies. (The idea of \textbf{confession} came from the church and is now used to collect data.)
\end{itemize}
\end{ExampleBox}

\begin{ExampleBox}
\textbf{Application of Foucault:}
\begin{itemize}
\item Social media (the IT Act, Section 66A); questions on \textbf{sexuality} (his book \textit{The History of Sexuality}); \textbf{surveillance}; the \textbf{census} and the \textbf{socio-economic caste census} (data collected but not published).
\end{itemize}
\end{ExampleBox}

\begin{QuickRevision}
\begin{itemize}
\item Power/knowledge/discourse intertwined; specialists manufacture our normal/abnormal; power is everywhere, not possessed.
\item Insanity/homosexuality were once normal --- modernity + psychology/sexology classified them as abnormal.
\end{itemize}
\end{QuickRevision}

\sectiondivider

\subsection{State-Centred Theory of Power (Theda Skocpol)}

\begin{KeyConcept}
\textbf{The state as autonomous:}
\begin{itemize}
\item \textbf{Theda Skocpol}: all previous theories (Marxist, feminist, pluralist, elite, power-elite, even functionalist) are \textbf{societal models} of power --- they project the state as a passive organ controlled by external agents (capitalists, men, interest groups, elites). She \textbf{rejects} them and proposes a \textbf{state-centred model}.
\item The state is an \textbf{autonomous institution} (with \textbf{relative autonomy}) that acts to secure \textbf{its own interest} --- not any class, section or gender. (Population control, surveillance --- the state acts for itself; even corporates and ex-IAS officers are surveilled.)
\end{itemize}
\end{KeyConcept}

\begin{PYQLink}
\begin{itemize}
\item 2019, Q4: ``What are the theoretical models of societal power? Which one of them is most applicable in advanced industrial society?'' --- here you write everything \textit{except} the state-centred theory (societal models = Marxist, feminist, pluralist, power-elite, functionalist).
\item No model is ``most'' applicable --- each is applicable in a context; show the \textit{application} of each (e.g., Marxist during lockdown when labour laws were suspended for FDI; feminist via Amartya Sen's ``missing women''/sarpanch-pati).
\end{itemize}
\end{PYQLink}

\begin{QuickRevision}
\begin{itemize}
\item Societal models (state = passive organ) vs Skocpol's state-centred model (state = autonomous, relative autonomy, acts for itself).
\item 2019 PYQ: list societal models; every theory is contextually applicable.
\end{itemize}
\end{QuickRevision}

\sectiondivider

\subsection{Applying the Power Elite to India (Yogendra Singh)}

\begin{DataFact}
\textbf{Power elite in Indian society:}
\begin{itemize}
\item The power elite (C. Wright Mills): heads of the three dominant institutions (military, corporate, federal executive) --- homogenous background, close group, \textbf{elite self-recruitment} $\rightarrow$ a \textbf{quasi-hereditary caste}. (The 9/11 attack on the World Trade Center, Pentagon and the federal government attacked all three symbols.)
\item \textbf{Yogendra Singh}: before independence, India's elites were \textbf{homogenous} (upper-caste, upper-class, Hindu men --- Brahmin-Kshatriya; Brahmins got land grants from the \textbf{Satavahana} kingdom to the \textbf{Gupta} period, and later to Buddhists). After independence (democracy, equal opportunity, land reforms, Zamindari Abolition, Land Ceiling Act, Cooperative Act), the elite structure became \textbf{heterogeneous} and fragmented --- Dalits (e.g., the President) and non-Hindus now occupy power.
\item \textbf{Importance of the power elite concept}: it formulates policies ``deemed fit'' (Mexican wall, visa curbs, dumping waste, war in the name of democracy); spreads democracy abroad (the \textbf{Arab Spring as a Facebook/American-propagation}); shows how democracy that \textit{appears} open is actually \textbf{closed} (needs of the few presented as needs of everyone --- e.g., anti-China propaganda); and \textbf{rejects the idea of democracy as an open society} (organized minority ruling unorganized masses).
\end{itemize}
\end{DataFact}

\begin{QuickRevision}
\begin{itemize}
\item Power elite = heads of military/corporate/federal; quasi-hereditary caste; 9/11 attacked the three symbols.
\item India: homogenous elite (pre-independence) $\rightarrow$ heterogeneous (post-independence, Yogendra Singh). Importance: policies ``deemed fit'', spread of democracy, democracy hides its own closure.
\end{itemize}
\end{QuickRevision}

\sectiondivider

\section{Civil Society \& Citizenship}

\subsection{Civil Society: Definition (Habermas \& the Public Sphere)}

\begin{KeyConcept}
\textbf{Definition:}
\begin{itemize}
\item \textbf{Civil society} = when \textbf{private individuals mobilise for a public cause} --- this is the \textbf{public sphere}. \textbf{Collective action is the spirit of civil society}. It lies \textbf{between the state and the family}.
\item The idea of the \textbf{public sphere} was coined by \textbf{J\"urgen Habermas} (mention him in any civil-society answer). The birth of sociology itself is the result of civil society (the Enlightenment); civil society led to the French Revolution, India's social-religious reforms and independence (Gandhi mobilised peasants, farmers, women, tribals).
\end{itemize}
\end{KeyConcept}

\begin{ExampleBox}
\begin{itemize}
\item The \textbf{anti-corruption movement}; the \textbf{Narmada Bachao Andolan} (informed citizens educating the masses about displacement); the \textbf{\#MeToo movement}; the \textbf{auto-rickshaw discussion} (deliberation over a policy) vs the \textbf{parliament debate} --- both are civil society. An \textbf{opposition party} can be civil society (a ruling party cannot); a private MP speaking for tribals (Shashi Tharoor) is civil society.
\end{itemize}
\end{ExampleBox}

\begin{QuickRevision}
\begin{itemize}
\item Civil society = private individuals mobilising for a public cause = public sphere (Habermas); collective action; between state and family.
\item Examples: anti-corruption, Narmada Bachao, \#MeToo, auto-rickshaw deliberation, opposition parties.
\end{itemize}
\end{QuickRevision}

\sectiondivider

\subsection{Features \& Types of Civil Society}

\begin{DataFact}
\textbf{Features:}
\begin{enumerate}
\item It creates \textbf{informed public opinion} on any policy/act of the state.
\item It \textbf{assists the state in formulating public policy}.
\item It \textbf{keeps check and balance on the power of the state}.
\end{enumerate}
\begin{itemize}
\item (World Economic Forum, UNESCO, UNCTAD, UNICEF, Bharatiya Kisan Union, student unions, opposition parties, media --- all civil society.)
\end{itemize}
\end{DataFact}

\begin{DataFact}
\textbf{Types:}
\begin{enumerate}
\item \textbf{Voluntary associations} (e.g., the JNU student union formed to protest a hostel rent hike, then dissolved).
\item \textbf{Faith/religious organisations} (Brahmo Samaj, Arya Samaj --- mobilising for reform of sati, child marriage, caste).
\item \textbf{NGOs, pressure groups, opposition parties, global institutions, media}.
\end{enumerate}
\end{DataFact}

\begin{KeyConcept}
\textbf{Neera Chandhoke:}
\begin{itemize}
\item \textbf{Neera Chandhoke}: civil society \textbf{neutralises individualism} in the wake of modernity and creates \textbf{solidarity/public cause across caste, class and religion} (e.g., \#BlackLivesMatter, \#DalitLivesMatter --- even non-Dalits/non-Christians express solidarity). (Political-administrative bureaucrats are part of the state, not civil society.)
\end{itemize}
\end{KeyConcept}

\begin{QuickRevision}
\begin{itemize}
\item Features: informed opinion, assist policy, check state. Types: voluntary, faith, NGO/pressure group/media/global.
\item Chandhoke: civil society neutralises individualism, creates solidarity across identities.
\end{itemize}
\end{QuickRevision}

\sectiondivider

\subsection{Roles of Civil Society}

\begin{DataFact}
\textbf{The four roles:}
\begin{enumerate}
\item \textbf{Policy advocacy} --- assisting the state to uplift weaker sections (UNESCO's \textbf{biosphere reserve} --- core/buffer zones --- protecting indigenous tribes; the \textbf{Biodiversity Act}; NGOs like Save the Children; women's/transgender associations giving inputs to legislation).
\item \textbf{Protection} --- providing \textbf{legal services/aid} to the marginalised (NALSA, National Legal Services Authority; \textbf{PIL --- public interest litigation}).
\item \textbf{Promotion of transparency} --- the \textbf{RTI campaign}, \textbf{Lokpal}, \textbf{Lokayukta} (all initiated by civil society, not the state).
\item \textbf{Mobilisation of citizens and resources} --- Swachh Bharat, Swachhta Pakhwada, plays, Bollywood awareness films.
\end{enumerate}
\end{DataFact}

\begin{CriticalPoint}
\textbf{Is civil society really ``civil''?:}
\begin{itemize}
\item Not always: NGOs may channelise foreign funds (FCRA --- Foreign Contribution Regulation Act) and appropriate funds for private interest; groups may serve their own ideology (RSS as ``uncivil society'' --- moral policing). \textbf{Partha Chatterjee} (Marxist): most Indian civil-society organisations serve the \textbf{elite}, not the poor --- with neoliberal/LPG reforms they have become ``liberalised''/market-oriented (landlords formed NGOs after Zamindari abolition to save land; more cow-vigilantism than transgender-vigilantism).
\end{itemize}
\end{CriticalPoint}

\begin{QuickRevision}
\begin{itemize}
\item Roles: policy advocacy (biosphere/Biodiversity Act), protection (NALSA, PIL), transparency (RTI/Lokpal), mobilisation (Swachh Bharat).
\item Uncivil society: FCRA funds, moral policing (RSS --- Chandhoke); Chatterjee: NGOs serve elites, not the poor.
\end{itemize}
\end{QuickRevision}

\sectiondivider

\subsection{Theories of Citizenship}

\begin{KeyConcept}
\textbf{Definition:}
\begin{itemize}
\item \textbf{Citizenship} = full participation of an individual in the social, economic and political sphere --- the \textbf{realisation of the rights of the individual}; it makes the state accountable. (Plato/Aristotle: citizenship is ancient --- Plato denied it to slaves, women, farmers; before the French Revolution the ``third estate'' bore the burden; the Enlightenment spoke of individuals' rights --- Mary Wollstonecraft extended citizenship to women.)
\end{itemize}
\end{KeyConcept}

\begin{DataFact}
\textbf{The theories:}
\begin{itemize}
\item \textbf{Marx}: citizenship is of no use if it does not create consciousness of exploitation --- proletarian citizenship cannot be exercised in a bourgeois society; societal transformation is inevitable. ``Class in itself'' is not citizenship; ``class for itself'' is.
\item \textbf{T. H. Green} (non-violent, same message): one cannot realise citizenship if one is not \textbf{well-placed in the social structures} (a minority, homosexual, poor, black --- victim of racism/casteism/homophobia despite being a citizen).
\item \textbf{Amartya Sen}: the \textbf{bundle of unfreedoms} (caste, gender, religion, ethnicity, sexuality, urban-rural divide); \textbf{development as freedom}; the \textbf{capability approach} (scholarships, women's education, girls' toilets, Kasturba Gandhi Balika Vidyalaya, Dalit capitalism/DICCI --- building people's capability, not just GDP).
\item \textbf{T. H. Marshall} (evolutionary theory of citizenship): British citizenship evolved in phases --- 18th century \textbf{civil citizenship} (civil rights: property, speech, justice); 19th century \textbf{political citizenship} (franchise, right to form associations); 20th century \textbf{social citizenship} (welfare schemes for women/black/poor). Unlike Britain, India granted citizenship ``at the stroke of midnight'' --- not phased.
\item (Note on \textbf{reservation}: education was the \textbf{monopoly of Brahmins for centuries}; today's reservation for weaker sections is a challenge to that age-old reservation --- hence its extension by ten years each time (the suffering of millions cannot be compensated in 10--20 years); \textbf{EWS} is a new, economically-based form of reservation.)
\end{itemize}
\end{DataFact}

\begin{QuickRevision}
\begin{itemize}
\item Citizenship = realisation of individual's rights. Marx (transformation inevitable), Green (social location), Sen (unfreedoms + capability), Marshall (civil $\rightarrow$ political $\rightarrow$ social, phased).
\end{itemize}
\end{QuickRevision}

\sectiondivider

\subsection{New Forms of Citizenship}

\begin{DataFact}
\begin{itemize}
\item \textbf{Environmental/ecological citizenship} --- the global cause of climate/environment now dominates discourse; it grants a right to a healthy life \textit{and} a duty to conserve nature (factory vs health --- the dialectic of developmentalism and environmentalism $\rightarrow$ \textbf{sustainable development}).
\item \textbf{Documentary citizenship} (Kamal Sadiq, \textit{Paper Citizens}) --- ``who has papers is a citizen'' (the Assam/NRC issue): the poor illegal migrant may get papers more easily (networks, wealth) than the poor indigenous citizen.
\item \textbf{Sexual citizenship} (Judith Butler, queer theory) --- citizenship rights are heteronormative; homosexuals are excluded.
\item \textbf{Dalit citizenship} --- Gail Omvedt, Sharmila Rege, Ambedkar (political freedom is useless without economic freedom; caste is ``division of labourers'', not just labour).
\item (Also: \textbf{political citizenship} lacking --- Baloch, Hong Kong, Kurds; \textbf{stateless nations} and \textbf{nationless states}.)
\end{itemize}
\end{DataFact}

\begin{QuickRevision}
\begin{itemize}
\item Environmental (right + duty), documentary (Paper Citizens --- Kamal Sadiq), sexual (Butler), Dalit (Omvedt/Rege/Ambedkar).
\end{itemize}
\end{QuickRevision}

\sectiondivider

\subsection{Nation, State, Nation-State \& State-Nation}

\begin{KeyConcept}
\textbf{The terms:}
\begin{itemize}
\item \textbf{Nation} = a group with similar culture, religion, belief and political consciousness. India is \textbf{not} a nation (multiple nations --- \textbf{Romila Thapar}); but India \textbf{is a state} (government + territory + sovereignty).
\item \textbf{Nation-state} = nation + state (unity in uniformity) --- a French idea (homogeneous society). \textbf{Yogendra Yadav} (founder of the Swaraj party; also in NCERT): India is a \textbf{state-nation} (the opposite) --- \textbf{unity in diversity}, not uniformity.
\item \textbf{Two states can become one} (Berlin Wall); \textbf{two nations can never become one}. A state may have multiple nations (India, UK, Spain/Basque, Canada/Quebec); a \textbf{stateless nation} (Baloch, Kurds, Gorkhaland); a \textbf{nationless state}; one nation in two states (Punjab of India and Pakistan); one state with two territories (USA + Alaska; East Pakistan $\rightarrow$ Bangladesh).
\end{itemize}
\end{KeyConcept}

\begin{QuickRevision}
\begin{itemize}
\item Nation (culture/belief) $\neq$ state (government+territory+sovereignty) $\neq$ nation-state (French, uniformity) vs state-nation (Yogendra Yadav, unity in diversity).
\item Stateless nation (Baloch/Kurds/Gorkhaland), nationless state, one nation two states (Punjab).
\end{itemize}
\end{QuickRevision}

\sectiondivider

\section{Ideology \& Pressure Groups}

\subsection{Ideology (Karl Mannheim: Ideology \& Utopia)}

\begin{KeyConcept}
\textbf{Definition:}
\begin{itemize}
\item \textbf{Ideology} = a set of ideas, beliefs or opinions. Every person is driven by an ideology; sociologists deal with it comprehensively (not just political ideologies like conservatism, liberalism, fascism, Nazism, neo-liberalism, neo-conservatism).
\item Everything we have studied is ideology: humanism (Renaissance), scientific revolution, Enlightenment, Reformation (which also produced sects --- highly absolutist, hence prone to \textbf{fundamentalism}, itself an ideology), French/industrial revolution, conservatism, structural functionalism, Marxism, positivism, feminism, nationalism, civil society. \textbf{Hegel's dialectical idealism} = the conflict of ideas.
\end{itemize}
\end{KeyConcept}

\begin{DataFact}
\textbf{Karl Mannheim (founder of the sociology of knowledge):}
\begin{itemize}
\item Book \textit{Ideology and Utopia}:
\begin{itemize}
\item \textbf{Ideology} = a set of ideas that promote the interest of the \textbf{ruling class} and maintain the \textbf{status quo} (e.g., functionalism --- ``every structure maintains the system'').
\item \textbf{Utopia} = a set of ideas that presents a \textbf{new form of world by challenging the existing world} (e.g., Marxism --- not ``absurd'', just a new/challenging idea).
\end{itemize}
\item Both functionalism and Marxism are \textbf{determinist} and \textbf{obscure the true nature of reality} --- all ideologies are \textbf{one-sided/one-dimensional} (feminism sees only patriarchy; Marxism sees only class).
\item \textbf{Marx}: every society is based on \textbf{ruling-class ideology / false consciousness} --- ideologies conceal reality. Every ideologue is a fundamentalist (left vs right = polar opposites of the same pole).
\end{itemize}
\end{DataFact}

\begin{CriticalPoint}
\textbf{Ideology unites and divides nation-states:}
\begin{itemize}
\item Nation-states are \textbf{formed} on the basis of ideologies (the French Revolution's homogeneous nationalism; the fall of the USSR $\rightarrow$ new nation-states on ethnic/religious/democratic lines) and \textbf{divided} by them (the India--Pakistan Partition = the clash of Gandhi's and Jinnah's ideologies). Ideologies also drive social movements (civil society, the Arab Spring).
\item \textbf{Science is itself an ideology} (objectivity, empiricism, replicability) --- recall the debate: \textbf{Popper} (religion closed, science open), \textbf{Kuhn} (both conservative and liberal), \textbf{Feyerabend} (``anything goes'').
\end{itemize}
\end{CriticalPoint}

\begin{QuickRevision}
\begin{itemize}
\item Ideology = set of ideas/beliefs; everything is ideology (humanism, science, revolution, functionalism, Marxism, feminism, nationalism).
\item Mannheim (Ideology \& Utopia): ideology = ruling-class/status-quoist (functionalism); utopia = new world (Marxism); both determinist and one-sided.
\item Nation-states formed/divided on ideology (USSR, Partition); science is an ideology (Popper/Kuhn/Feyerabend).
\end{itemize}
\end{QuickRevision}

\sectiondivider

\subsection{Pressure Group: Definition \& Functions}

\begin{KeyConcept}
\textbf{Definition:}
\begin{itemize}
\item A \textbf{pressure group} = a set of people mobilised for a cause \textbf{beyond the system of party politics} --- it operates outside party politics and mobilises individuals for a specific (sectional) cause. There is \textbf{no collective/national interest} --- only sectional interests. (The idea is pluralist, from Robert Dahl --- the first pressure groups were conceptualised in America.)
\item A pressure group is a \textbf{theoretical model} --- nobody calls themselves a ``pressure group''; they call themselves associations.
\end{itemize}
\end{KeyConcept}

\begin{DataFact}
\textbf{Functions:}
\begin{enumerate}
\item \textbf{Represents sectional interests} (e.g., for/against marital rights for transgenders).
\item \textbf{Political socialisation} (informs people of their rights and mobilises them --- candle marches, pamphlets).
\item \textbf{Checks the authoritarian tendency of the state} (Indira Gandhi's authoritarian regime $\leftrightarrow$ J. P. Narayan's movement; Rajni Kothari's \textbf{``one-party dominance''/Congress system} --- now the new one-party dominance of the BJP).
\item \textbf{Influences legislation} (VHP, RSS, FICCI, ASSOCHAM, Bengal Chamber of Commerce, student unions, AITUC, Bharatiya Kisan Union, PETA, Amnesty International, Pinjra Tod).
\end{enumerate}
\end{DataFact}

\begin{QuickRevision}
\begin{itemize}
\item Pressure group = mobilisation beyond party politics for a sectional cause (pluralist, Dahl).
\item Functions: represent sectional interest, political socialisation, check authoritarianism (Kothari's one-party dominance), influence legislation.
\end{itemize}
\end{QuickRevision}

\sectiondivider

\subsection{Classification of Pressure Groups (Almond \& Powell)}

\begin{DataFact}
\textbf{Four types (Almond \& Powell):}
\begin{enumerate}
\item \textbf{Associational} --- registered bodies pursuing limited goals for certain sections (AITUC for workers; DUTA; civil-services officers' institute).
\item \textbf{Non-associational} --- informal, may or may not be registered (caste groups --- Harijan Sevak Sangh, Marwari Association, Bhim Army; ideology groups --- Narmada Bachao; tribal groups --- NSCN, Jharkhand Mukti Morcha).
\item \textbf{Institutional} --- pressure groups within the state apparatus (IAS/IPS associations, Army Wives Welfare Association, Air Force Wives Welfare Association).
\item \textbf{Anomic} --- resort to illegitimate/violent means (Naxalites, secessionists, Kashmir separatists/Hurriyat, reservation vandals, NSCN-Khaplang).
\end{enumerate}
\end{DataFact}

\begin{KeyConcept}
\textbf{Maurice Duverger:}
\begin{itemize}
\item \textbf{Exclusive pressure groups} (fully committed to challenging the state, nothing else) vs \textbf{partial pressure groups} (challenge the state for specific issues, then return to other work). \textbf{Pseudo pressure groups} work for a third party for monetary gain (media channels --- Duverger's own example).
\end{itemize}
\end{KeyConcept}

\begin{CriticalPoint}
\textbf{Pressure groups in India (Myron Weiner):}
\begin{itemize}
\item \textbf{Myron Weiner}: the Indian state \textbf{exploits pressure groups on the basis of numbers/vote banks} --- it hears VHP/RSS (numbers = votes) more than NGOs like the Naas Foundation (the exploited, without numbers). Pressure groups in India have been turned upside-down --- only those that can empower the state are heard; they ``play hide-and-seek in politics'' (is it resistance, or a route to power?).
\end{itemize}
\end{CriticalPoint}

\begin{QuickRevision}
\begin{itemize}
\item Almond \& Powell: associational, non-associational, institutional, anomic.
\item Duverger: exclusive / partial / pseudo. Weiner: state exploits pressure groups as vote banks.
\end{itemize}
\end{QuickRevision}

\sectiondivider

\section{Social Stratification --- Structural Functional Theory}

\subsection{Answer-Writing: Ideology \& Social Transformation}

\begin{PYQLink}
\begin{itemize}
\item ``Ideology is crucial for social transformation in a democracy.'' (10 marks.) Approach: for 10-mark questions, show \textbf{multiple dimensions with minimum explanation}; for 20-mark, limited dimensions with full explanation. Keywords: ideology, social transformation, democracy.
\end{itemize}
\end{PYQLink}

\begin{KeyConcept}
\textbf{The model:}
\begin{itemize}
\item Any society begins \textbf{relatively closed} (not democratic/inclusive); an \textbf{ideology} comes in, and through the reaction of ideology + closed society it becomes \textbf{relatively more open} (+ byproducts, e.g., fundamentalism). A new ideology makes it even more open --- this is \textbf{social transformation}.
\end{itemize}
\end{KeyConcept}

\begin{ExampleBox}
\textbf{India:}
\begin{itemize}
\item Pre-independence India was closed (only Rajas/Maharajas in legislative assemblies). The \textbf{exogenic} ideology of \textbf{liberty, equality, fraternity, justice} (borrowed from the West, enshrined in the Preamble) made it open (meritocracy). But it was closed at the top (Congress dominance till the mid-70s --- C. Wright Mills' ``external appearance of democracy''). Then \textbf{regionalism} (an orthogenic ideology --- Reddys, Jats, Chaudharis, Gujjars; the Mandal Commission 1979, VP Singh, OBC reservation affirmed by the Supreme Court in 1992; coalition politics) made it more open. Then \textbf{LPG/neo-liberalism} (1991) opened the economy, and with globalisation came \textbf{feminism, democratic decentralisation (73rd/74th Amendment, Panchayati Raj), environmentalism, RTI/transparency, individualism} --- further social transformation.
\end{itemize}
\end{ExampleBox}

\begin{QuickRevision}
\begin{itemize}
\item Model: closed society + ideology $\rightarrow$ more open (social transformation). 10-mark = multiple dimensions, minimal explanation.
\item India: exogenic (liberty/equality) $\rightarrow$ regionalism/Mandal (orthogenic) $\rightarrow$ LPG + feminism/decentralisation (1991 onwards).
\end{itemize}
\end{QuickRevision}

\sectiondivider

\subsection{Social Stratification: An Introduction}

\begin{KeyConcept}
\begin{itemize}
\item \textbf{Stratification} = layers of people, just like the strata of the earth. Society is divided into strata everywhere (Chanakyapuri embassies vs Old Rajinder Nagar; coaching payers vs Telegram-users). A strata can be hierarchical.
\item These strata are \textbf{models/ideal types} (Weber) --- not reality. Class-based strata (lower/middle/upper class + elites) are \textbf{open} (an auto-driver's son can become IAS; Dhirubhai Ambani to Mukesh Ambani). \textbf{Varna} (Shudra--Vaishya--Kshatriya--Brahmin) was \textit{philosophically open} (Gandhi: ``we are all born Shudras, but by talent become Brahmin''; the Triguna theory --- Rajasvik/Sattvik/Tamasik qualities determine varna), while \textbf{caste (jati)} is \textbf{closed} --- varna unites, caste divides.
\item Beyond the fourfold varna are the \textbf{Pancham Varna} (outcasts/untouchables), the \textbf{Mleccha} (foreigners), the \textbf{Chandal}, and the \textbf{tribals --- who were never part of the Hindu system}.
\item Stratification is an \textbf{unconscious device} --- nobody planned it; it exists in every society.
\end{itemize}
\end{KeyConcept}

\begin{QuickRevision}
\begin{itemize}
\item Strata = layers (models/ideal types). Class = open; varna = philosophically open (Triguna); caste (jati) = closed; varna unites, caste divides.
\item Stratification = unconscious, inevitable, organic.
\end{itemize}
\end{QuickRevision}

\sectiondivider

\subsection{Structural Functional Theory (Parsons, Davis \& Moore)}

\begin{KeyConcept}
\textbf{Parsons:}
\begin{itemize}
\item \textbf{Talcott Parsons} (American, ``land of opportunity''): strata are based on \textbf{class position}, a function of individual skill and merit. The \textbf{school is a focal agency} between family and society --- it transmits the values (meritocracy, equality of opportunity) that form \textbf{value consensus}, so those who don't succeed blame themselves, not the system. Some positions are more important than others; \textbf{school does effective role allocation} (monitor, captain) --- extrapolated to the wider society; school = a \textbf{mini society}.
\end{itemize}
\end{KeyConcept}

\begin{DataFact}
\textbf{Davis \& Moore (structural functional analysis of stratification):}
\begin{itemize}
\item Some positions are functionally more important; society must attract the most meritorious/talented to them by attaching \textbf{rewards} (perks, allowances, pension, prestige --- e.g., the IAS). People compete on \textbf{merit}; those not selected accept defeat (system is merit-based). So \textbf{inequality is functional} --- it recruits meritorious people; conflict is neutralised; the system is maintained. Stratification is an \textbf{unconscious device}.
\end{itemize}
\end{DataFact}

\begin{QuickRevision}
\begin{itemize}
\item Parsons: meritocracy, school = focal agency, effective role allocation, value consensus.
\item Davis \& Moore: important positions + rewards $\rightarrow$ meritorious recruitment; inequality is functional; system maintained.
\end{itemize}
\end{QuickRevision}

\sectiondivider

\subsection{Melvin Tumin's Critique}

\begin{CriticalPoint}
\textbf{Tumin's four objections:}
\begin{enumerate}
\item \textbf{How do you decide which position is ``important''?} (Is an IAS officer more important than a Bollywood actor? Why is the IAS paid $\sim$\rupee 60,000 while the actor earns crores? Status inconsistency --- Gerhard Lenski.)
\item There is \textbf{no objectivity} --- importance is \textbf{spatial and temporal} (doctors became ``medical warriors'' during COVID; ASHA/Anganwadi workers' emoluments raised). Importance = shared agreement, not objective fact.
\item \textbf{Not everyone is given equal opportunity} --- transgenders couldn't appear for UPSC until 2018--19; women, Dalits, the poor are excluded.
\item If the system is merit-based, \textbf{why are important positions always occupied by people from a specific social background} (white, upper-class, Ivy League, Protestant --- in India, upper-caste men)? Davis \& Moore \textbf{ignore the social location} (race, caste, gender, ethnicity) of individuals.
\end{enumerate}
\begin{itemize}
\item Hence stratification is \textbf{dysfunctional}, not functional. (Also: Davis \& Moore cannot explain the uneducated social-media influencers with wealth/status/power --- social media is the new focal agency.)
\end{itemize}
\end{CriticalPoint}

\begin{QuickRevision}
\begin{itemize}
\item Tumin: (1) how to judge ``importance'', (2) importance is spatial/temporal, (3) no equal opportunity, (4) important positions held by specific social backgrounds.
\item Stratification = dysfunctional; Davis \& Moore ignore social location.
\end{itemize}
\end{QuickRevision}

\sectiondivider

\section{Marxist \& Weberian Theories of Stratification}

\subsection{Marxist Theory of Stratification}

\begin{KeyConcept}
\begin{itemize}
\item To Marx the world has \textbf{two classes}: owners of the means of production (bourgeoisie/ruling class) and those who sell their labour (proletariat/exploited class). Strata = one's \textbf{relation to the means of production}. Power and status are \textbf{derived from the economy} (monocausal).
\item The working class will develop consciousness, overthrow the bourgeoisie (\textbf{proletarian revolution}), and establish a \textbf{classless society} (inequality $\rightarrow$ equality). \textbf{Polarisation} of classes is inevitable (the middle class will shrink --- polarise into bourgeoisie/proletariat). \textbf{Collective action} of the proletariat is inevitable.
\end{itemize}
\end{KeyConcept}

\begin{QuickRevision}
\begin{itemize}
\item Marx: two classes (bourgeoisie/proletariat); strata = relation to means of production; polarisation + revolution $\rightarrow$ classless society; collective action inevitable.
\end{itemize}
\end{QuickRevision}

\sectiondivider

\subsection{Weberian Theory of Stratification (Class, Status, Party)}

\begin{KeyConcept}
\textbf{Class:}
\begin{itemize}
\item \textbf{Class} = a group of people with a \textbf{similar market situation} by virtue of which they enjoy \textbf{similar life chances}. (An engineer and a professor both earn $\sim$70,000 $\rightarrow$ similar market situation and life chances; but a cobbler, though skilled, has a different market position, no prestige/power.) Market situation is judged by \textbf{wealth and income} (not just ownership of the means of production).
\item Weber \textbf{agrees} with Marx's division by means of production but says it is \textbf{not the only} division. He rejects monocausality --- \textbf{status and power can also lead to wealth} (multi-causal/adequate causality).
\end{itemize}
\end{KeyConcept}

\begin{KeyConcept}
\textbf{Status group \& party:}
\begin{itemize}
\item \textbf{Status group} = people with similar \textbf{esteem/status/prestige} (Scindia, Gaekwad, the royal family of Bikaner --- different wealth, same status). Status governs \textbf{lifestyle}; class governs \textbf{life chances}. Wealth and status have \textbf{no necessary link} --- a \textbf{Dawood Ibrahim} has wealth but no status, while a bankrupt \textbf{Vijay Mallya} still commands status/honour.
\item \textbf{Party} = people with similar \textbf{interests and power} (a pressure group, a Naxalite party, a trade union, moderates/extremists --- \textit{not} necessarily a political party). The working class is \textbf{not homogeneous} --- it forms different parties (hence no inevitable collective action).
\item Weber's \textbf{trinitarian model}: stratification on the basis of \textbf{wealth (class), status, and power} --- three separate continua with \textbf{no necessary relationship} between them. (Marx: power/status derive from economy; Weber: the reverse is also possible.)
\end{itemize}
\end{KeyConcept}

\begin{QuickRevision}
\begin{itemize}
\item Class = market situation + life chances (wealth + income). Status group = esteem/prestige (lifestyle). Party = similar interests/power.
\item Trinitarian model; multi-causal (not monocausal); working class is heterogeneous (different parties).
\end{itemize}
\end{QuickRevision}

\sectiondivider

\subsection{Weber's Ideal Types of Class}

\begin{DataFact}
\textbf{The four ideal types:}
\begin{enumerate}
\item \textbf{Propertied upper class} (bourgeoisie --- own factory/land).
\item \textbf{Property-less white-collar workers} (professors, engineers, lawyers, doctors, pilots --- they may have a house, but not the means of production).
\item \textbf{Petty bourgeoisie} (small shopkeepers, micro/small industries, startups --- employers of fewer than $\sim$10).
\item \textbf{Menial labourers} (auto/rickshaw drivers, dhobi, mochi, carpenter).
\end{enumerate}
\begin{itemize}
\item Marx would collapse these into bourgeoisie + proletariat; Weber \textbf{introduces the middle class} and says it will \textbf{expand/inflate} with rationalisation/specialisation (vs Marx's \textbf{polarisation} --- the middle class shrinking). Weber: \textbf{diversification of class}; Marx: polarisation.
\end{itemize}
\end{DataFact}

\begin{QuickRevision}
\begin{itemize}
\item Four ideal types: propertied upper, property-less white-collar, petty bourgeoisie, menial labourers.
\item Weber: middle class expands (rationalisation/specialisation); Marx: polarises/shrinks.
\end{itemize}
\end{QuickRevision}

\sectiondivider

\subsection{Erik Olin Wright: Contradictory Class Location}

\begin{KeyConcept}
\begin{itemize}
\item \textbf{Erik Olin Wright} (1960s--70s) rejects both Marx and Weber: classes are not homogeneous --- people occupy \textbf{contradictory class locations}. A manager (white-collar) \textit{controls} workers like a bourgeois, but is \textit{himself controlled} by the bourgeois. In finance capitalism, one may be a worker in Amazon \textit{and} a shareholder of Reliance. He classifies by \textbf{control over workers and ownership of the means of production} (both in degree): upper class (owns + controls, employs managers), middle class (controls partially + employs working class), working class (neither owns nor can recruit --- sells labour).
\end{itemize}
\end{KeyConcept}

\begin{QuickRevision}
\begin{itemize}
\item Olin Wright: contradictory class location (manager controls workers but is controlled); finance capitalism (worker + shareholder); classes by control + ownership.
\end{itemize}
\end{QuickRevision}

\sectiondivider

\subsection{Frank Parkin: Social Closure}

\begin{KeyConcept}
\textbf{Social closure:}
\begin{itemize}
\item \textbf{Frank Parkin} (Marxist \textit{and} Weberian): classes should be defined by the \textbf{strategy employed to monopolise resources} --- \textbf{social closure}.
\item \textbf{Exclusionary social closure} --- insiders exclude outsiders (the UPSC's English and Central-India-focused history; the air-force simulation test favouring commanders' children; St. Stephen's/Hindu/DSE cut-offs and interviews; recruitment via insider contacts --- ``opportunity hoarding''; insider trading; elite self-recruitment). It is the strategy of the \textbf{dominant class}.
\item \textbf{Usurpationary social closure} --- the excluded class unites to \textbf{claim back} hoarded resources (feminism vs male-monopolised occupations; peasant movements vs landlords). It is the strategy of the \textbf{oppressed class}.
\item \textbf{Dual closure} --- once the excluded gain opportunity, they themselves exclude others (Dalit elites exclude Dalit women; Mayawati celebrated \textit{because} she is both Dalit \textit{and} a woman). The \textbf{middle class} exercises dual closure.
\item (Examples: \textbf{Kancha Ilaiah} on Dalit women; the \textbf{Nirbhaya vs Hathras} difference --- upper-caste victim vs Dalit victim, treated differently.)
\end{itemize}
\end{KeyConcept}

\begin{QuickRevision}
\begin{itemize}
\item Parkin: class = strategy of resource monopolisation. Exclusionary (dominant class), usurpationary (oppressed), dual (middle class).
\item Nirbhaya (upper-caste) vs Hathras (Dalit woman) --- unequal treatment.
\end{itemize}
\end{QuickRevision}

\sectiondivider

\subsection{The ``Death of Class'' Thesis (Waters \& Pakulski)}

\begin{KeyConcept}
\textbf{Post-class society:}
\begin{itemize}
\item \textbf{Malcolm Waters \& Jan Pakulski} (\textit{Death of Class}): class is no longer an important criterion. With welfare states, equal opportunity, globalisation/liberalisation and democratic revolutions, wealth has been distributed, class differences have been subjugated, and the working class no longer votes on class lines (e.g., workers not voting for the Communist Party).
\item \textbf{Post-modern politics} is based on \textbf{non-economic issues --- ethnicity, identity, environment}. Society has shifted from an \textbf{economic class society to a status society} (lifestyle --- what you wear, eat, shop --- defines you, not class).
\end{itemize}
\end{KeyConcept}

\begin{QuickRevision}
\begin{itemize}
\item Waters \& Pakulski: death of class; post-modern politics on ethnicity/identity/environment; class society $\rightarrow$ status society (lifestyle).
\end{itemize}
\end{QuickRevision}
